% ============================================================
% IEEE TEMSMET 2026 — FINAL SUBMISSION VERSION
% Upload as main.tex at Overleaf project root.
% Figures: outputs/figures/ subfolder of project root.
% Compile: pdflatex -> bibtex -> pdflatex -> pdflatex
%
% NOTE FOR AUTHOR (optional enhancement):
% To add an Experimental Workflow diagram as Figure 1, export
% the pipeline SVG to outputs/figures/experimental_workflow.png
% (300 DPI, A4 width) and uncomment the figure block below the
% \maketitle command. All subsequent figures will renumber
% automatically.
% ============================================================

\documentclass[conference,a4paper]{IEEEtran}

% --- Packages ---------------------------------------------
\usepackage[T1]{fontenc}
\usepackage[utf8]{inputenc}
\usepackage{amsmath,amssymb}
\usepackage{bm}
\usepackage{graphicx}
\usepackage{booktabs}
\usepackage{array}
\usepackage{url}
\usepackage{xspace}
\usepackage{cite}
\usepackage{microtype}
\usepackage[hidelinks]{hyperref}

% --- Figure path (Overleaf: main.tex at project root) -----
\graphicspath{{outputs/figures/}}

% --- Macros -----------------------------------------------
\newcommand{\SHAP}{\textsc{shap}\xspace}
\newcommand{\LIME}{\textsc{lime}\xspace}
\newcommand{\NIDS}{NIDS\xspace}
\newcommand{\SMOTE}{SMOTE\xspace}
\newcommand{\XAI}{XAI\xspace}
\newcommand{\rbs}{rank-biserial~$r$\xspace}
\newcommand{\cohensh}{Cohen's~$h$\xspace}

% ==========================================================
\begin{document}
% ==========================================================

\title{Impact of SMOTE-Based Class Rebalancing\\
       on SHAP and LIME Explanation Quality\\
       in Random Forest Network Intrusion Detection}

\author{%
  \IEEEauthorblockN{Yogesh Narendran}
  \IEEEauthorblockA{%
    Department of Computer Science Engineering\\
    (Artificial Intelligence and Machine Learning)\\
    Sri Sairam Engineering College\\
    Chennai, Tamil Nadu, India\\
    Email: \texttt{sec24am081@sairamtap.edu.in}}}

\maketitle

% ----------------------------------------------------------
\begin{abstract}
Machine learning--based Network Intrusion Detection Systems (\NIDS) routinely
apply the Synthetic Minority Over-sampling Technique (\SMOTE) to address severe
traffic-class imbalance, yet the impact of such rebalancing on post-hoc
explainability outputs has not been systematically quantified. This study
investigates how \SMOTE rebalancing affects feature attribution rankings produced
by SHapley Additive exPlanations (\SHAP) and Local Interpretable Model-agnostic
Explanations (\LIME) in a Random Forest \NIDS trained and evaluated on the
UNSW-NB15 benchmark. Two classifiers --- one trained on the original imbalanced
partition and one on a \SMOTE-rebalanced partition --- are evaluated against the
same 82,332-instance test set using a statistical framework comprising McNemar's
test, Wilcoxon signed-rank tests, bootstrap confidence intervals, and
Holm--Bonferroni correction. \SMOTE improves macro-averaged F1 and substantially
increases minority-class recall, with all aggregate predictive effect sizes
remaining negligible (\cohensh $\leq 0.09$). In contrast, \SHAP attribution
rankings shift with a medium effect ($r=0.33$) and \LIME rankings with a large
effect ($r=0.60$), demonstrating that class rebalancing alters explanation outputs
disproportionately relative to classification metrics. Accuracy-only validation
after \NIDS retraining is insufficient to characterise post-hoc explanation
stability.
\end{abstract}

\begin{IEEEkeywords}
network intrusion detection, explainability, class imbalance,
synthetic minority oversampling, Shapley additive explanations,
local interpretable model-agnostic explanations, random forest,
feature attribution
\end{IEEEkeywords}

% ==========================================================
\section{Introduction}
% ==========================================================

Machine learning--based \NIDS have demonstrated strong performance on curated
benchmarks \cite{alshamy2021,wu2022,more2024}, yet a persistent challenge is
severe class imbalance: benign traffic dominates by orders of magnitude while rare
attack categories may constitute fewer than 0.1\% of packets
\cite{moustafa2015,shanmugam2025}. This asymmetry biases classifiers toward
majority-class predictions and suppresses detection of rare, high-severity attacks.

\SMOTE \cite{chawla2002} is the dominant remediation strategy, generating
synthetic minority-class instances by interpolating between existing samples in
feature space. \SMOTE has been shown to improve minority-class recall in \NIDS
applications \cite{alshamy2021,sayegh2024}; in this study it increased Worms-class
F1 from 0.233 to 0.452.

Concurrently, eXplainable AI (\XAI) adoption in security operations has grown
substantially \cite{charmet2022,rjoub2023}. As AI systems become subject to
regulatory transparency requirements \cite{euaiact2024}, the reliability of
post-hoc explanations across model updates is critical. \SHAP
\cite{lundberg2017,lundberg2020} and \LIME
\cite{ribeiro2016,gaspar2024,hermosilla2025a} are the two most widely applied
post-hoc explanation methods for \NIDS models.

A critical and underexplored question is: \emph{when a \NIDS model is retrained
with a different class distribution, does the explanation change as much as the
prediction?} If explanation instability is disproportionate to predictive
instability, operators relying on accuracy alone to validate retrained models may
unknowingly deploy systems whose explanation outputs have changed substantially,
undermining the trustworthiness of \XAI-assisted security analysis.

This paper addresses that gap with a controlled empirical study. The contributions
are:
\begin{enumerate}
  \item A paired statistical comparison (McNemar's test, Wilcoxon signed-rank,
        95\% bootstrap CI) of Baseline and \SMOTE-trained Random Forest models
        on UNSW-NB15 (82,332 test instances).
  \item Quantified effect sizes for predictive metrics (\cohensh) and explanation
        stability (\rbs), with explicit effect-size interpretability bounds.
  \item A systematic comparison of \SHAP and \LIME sensitivity to training class
        distribution, demonstrating that \LIME is substantially more sensitive.
  \item Actionable guidance on explanation-consistency monitoring in
        \XAI-augmented \NIDS pipelines.
\end{enumerate}

% ==========================================================
\section{Related Work}
\label{sec:related}
% ==========================================================

\subsection{Class Imbalance in NIDS}

Class imbalance is a well-documented challenge in \NIDS
\cite{moustafa2015,shanmugam2025}. Minority attack categories are rare relative
to normal traffic, causing classifiers to achieve high accuracy by predicting the
majority class \cite{alshamy2021,chawla2002}. Re-sampling strategies ---
including \SMOTE, ADASYN, and random undersampling \cite{shanmugam2025} --- along
with cost-sensitive learning and ensemble methods \cite{shanmugam2025} have been
proposed. \SMOTE remains dominant due to its consistency and
reproducibility \cite{wu2022,sayegh2024}. UNSW-NB15 \cite{moustafa2015} is the
established benchmark used here, selected for its predefined train/test split and
role as a community evaluation standard \cite{more2024,hermosilla2025a}.

\subsection{Explainable AI for Network Security}

\XAI application to \NIDS has expanded alongside the trustworthy-AI
movement \cite{charmet2022,rjoub2023}. \SHAP provides attributions with
theoretical guarantees of consistency and local
accuracy \cite{lundberg2017,lundberg2020}. \LIME approximates model behaviour
locally with an interpretable surrogate without requiring structural knowledge of
the model \cite{ribeiro2016}. Both have been applied to Random
Forest \cite{hermosilla2025a,hermosilla2025b} and deep learning--based
\NIDS \cite{gaspar2024}.

\subsection{XAI Stability and Explanation Quality}

Explanation stability under perturbation has been studied in general ML
contexts \cite{visani2022} but remains underexplored in \NIDS settings. Prior
work examined adversarial perturbation effects on \LIME and \SHAP
\cite{charmet2022} and sensitivity to hyperparameters \cite{visani2022}. The
specific question of how training data rebalancing affects explanation stability
--- distinct from prediction stability --- has not been addressed in prior \NIDS
literature \cite{patil2020,alarab2022}. This study addresses that gap directly.

% ==========================================================
\section{Methodology}
\label{sec:method}
% ==========================================================

\subsection{Dataset}

\begin{figure}[t]
  \centering
  \includegraphics[width=\columnwidth]{class_distribution.png}
  \caption{Training set class distribution before \SMOTE (175,341 instances,
           10 classes, log scale). Maximum imbalance: 430.77:1
           (Worms: 130 vs.\ Normal: 56,000). After \SMOTE: 56,000 per class.}
  \label{fig:class_dist}
\end{figure}

The UNSW-NB15 dataset \cite{moustafa2015} provides a predefined training set
(175,341 instances) and test set (82,332 instances) across 10 classes. The
training set exhibits severe imbalance (Fig.~\ref{fig:class_dist}): Normal
contains 56,000 samples (31.9\%) while Worms contains only 130 (0.074\%),
yielding a maximum imbalance ratio of 430.77:1. UNSW-NB15 was selected for its
controlled-environment realism, predefined split, and status as a community
benchmark for reproducible \NIDS evaluation \cite{more2024}.

\subsection{Feature Engineering}

The original 44-feature schema was reduced to 42 features by removing
\texttt{id} (identifier) and \texttt{label} (target leakage). No scaling was
applied; Random Forest is scale-invariant.

\subsection{Class Rebalancing}

\SMOTE \cite{chawla2002} was applied with \texttt{sampling\_strategy=`auto'},
\texttt{k\_neighbors=5}, seed~42. Minority classes were oversampled to 56,000
per class, producing 560,000 training rows (384,659 synthetic samples). \SMOTE
was selected as the most widely benchmarked \NIDS rebalancing technique
\cite{alshamy2021,wu2022}, enabling direct comparison with prior work. The test
set was not modified; both models are evaluated on the same held-out set.

\subsection{Classifier}

A Random Forest \cite{breiman2001} --- selected for strong \NIDS benchmark
performance and native compatibility with \SHAP TreeExplainer --- was trained
independently on the Baseline (imbalanced) and \SMOTE-balanced datasets. All
hyperparameters and seed~(42) were held constant. The original train/test split
was used without re-splitting.

\subsection{XAI Methods}

\noindent\textbf{SHAP:} \SHAP values \cite{lundberg2017,lundberg2020} were
computed for 60 stratified test instances (six per class) using
\texttt{shap.TreeExplainer}, yielding exact Shapley values. Global importance
was obtained by averaging $|\text{SHAP}|$ values. Rankings were compared using
Spearman~$\rho$ and top-5 overlap.

\smallskip\noindent\textbf{LIME:} \LIME explanations \cite{ribeiro2016} were
computed on the same 60 instances using \texttt{LimeTabularExplainer}
(\texttt{num\_features=10}, \texttt{num\_samples=5000}). Surrogate fidelity
was measured by local $R^2$. The deterministic shared-instance protocol
eliminates selection variance from the paired comparison.

\subsection{Statistical Validation}
\label{sec:stats}

Four hypothesis tests (RQ1--RQ4) were conducted (Table~\ref{tab:hyp}):
(1)~\textbf{McNemar's test} is the appropriate paired test for comparing two
classifiers on the same test set. (Yates correction) on the 82,332-sample
prediction matrix; discordant pairs: $b=1{,}632$ (SMOTE gains), $c=4{,}784$
(SMOTE losses).
(2)~\textbf{Wilcoxon signed-rank} on paired confidence scores (82,332 pairs,
34,002 pairs with zero difference were excluded using the Wilcoxon zero-method). (3)~\textbf{Wilcoxon} on 42 paired \SHAP
importances. (4)~\textbf{Wilcoxon} on 36 paired \LIME importances.

Holm--Bonferroni correction was applied ($\alpha/4,\,\alpha/3,\,\alpha/2,\,\alpha$).
Effect sizes use two complementary metrics. \cohensh quantifies the difference
between two proportions:
\begin{equation}
  h = 2\arcsin\!\sqrt{p_1} - 2\arcsin\!\sqrt{p_2},
  \label{eq:cohenh}
\end{equation}
where $|h|<0.20$ is negligible, $\geq0.20$ small, $\geq0.50$ medium, $\geq0.80$
large. \rbs is derived from the Wilcoxon statistic:
\begin{equation}
  r = \frac{Z}{\sqrt{N}},
  \label{eq:rbs}
\end{equation}
where $N$ is non-zero paired differences; $r\geq0.10$ small, $\geq0.30$ medium,
$\geq0.50$ large. Because \cohensh operates on proportion data ($n = 82{,}332$)
and \rbs operates on feature importance ranks ($n = 42$ or $36$), cross-metric
comparisons are therefore interpreted as qualitative indicators. Power: \SHAP
at $n=42$, $r=0.326$ gives $\approx0.55$; \LIME at $n=36$, $r=0.595$ gives
$\approx0.90$. Bootstrap CIs: 95\%, 2,000 resamples, percentile method.

% ==========================================================
\section{Experimental Setup}
\label{sec:setup}
% ==========================================================

All experiments were implemented in Python~3.10+ using scikit-learn
\cite{pedregosa2011}, imbalanced-learn \cite{lemaitre2017}, \texttt{shap}
\cite{lundberg2017}, \texttt{lime} \cite{ribeiro2016}, and SciPy
\cite{virtanen2020}. The global random seed was fixed at~42. All configuration
parameters are stored in \texttt{configs/*.yaml}; all output artefacts are
SHA-256 verified before statistical analysis.

% ==========================================================
\section{Results}
\label{sec:results}
% ==========================================================

\subsection{Predictive Performance (RQ1)}

\begin{figure}[t]
  \centering
  \includegraphics[width=0.49\columnwidth]{confusion_matrix_baseline.png}\hfill
  \includegraphics[width=0.49\columnwidth]{confusion_matrix_smote.png}
  \caption{Normalised confusion matrices. (left)~Baseline; (right)~\SMOTE.
           \SMOTE improves minority-class recall (Worms, Analysis, Backdoor)
           at a marginal cost to majority-class precision.}
  \label{fig:cm}
\end{figure}

Table~\ref{tab:perclass} reports per-class metrics; Table~\ref{tab:bootstrap}
reports 95\% bootstrap confidence intervals.

\SMOTE reduced accuracy by 0.0383 ($0.7543\to0.7161$; 95\%~CI Baseline
$[0.7516,\,0.7574]$ vs.\ SMOTE $[0.7129,\,0.7191]$; non-overlapping) and
weighted F1 by 0.0149, while improving macro F1 by 0.0170. McNemar's test
confirmed statistically significant model disagreement
($\chi^2_{\text{Yates}}(1)=1547.51$, $p\approx 0$; corrected $\alpha=0.0125$).
Despite statistical significance, all \cohensh values are negligible:
$h(\text{accuracy})=-0.0868$, $h(\text{macro~F1})=+0.0340$,
$h(\text{weighted~F1})=-0.0354$ (Table~\ref{tab:effect}).

At the class level (Fig.~\ref{fig:cm}, Table~\ref{tab:perclass}), \SMOTE
improved F1 for minority attack classes (mean recall gain $+0.325$; Worms:
$0.233\to0.452$) while degrading F1 for majority classes (mean F1 change
$-0.020$; Normal: $0.848\to0.809$). The Worms test class contains only
44~instances; per-class F1 statistics at this sample size are sensitive to
individual prediction outcomes and are presented as indicative results rather
than stable estimates of population-level performance.

\smallskip
\noindent\textbf{Summary (RQ1):} Statistically significant but practically
negligible aggregate differences; minority-class F1 improved substantially while
majority-class F1 declined modestly.

% --- TABLE I (full-width) ----------------------------------
\begin{table*}[t]
  \caption{Per-Class Classification Metrics --- Baseline and SMOTE Models
           (Test Set, $n=82{,}332$)}
  \label{tab:perclass}
  \centering
  \begin{tabular}{lrrrrrrrr}
    \toprule
    \textbf{Class} &
      \textbf{P\,(B)} & \textbf{P\,(S)} &
      \textbf{R\,(B)} & \textbf{R\,(S)} &
      \textbf{F1\,(B)} & \textbf{F1\,(S)} &
      $\bm{n}$ \\
    \midrule
    Analysis          & 0.005 & 0.034 & 0.009 & 0.164 & 0.006 & 0.057 & 677 \\
    Backdoor          & 0.017 & 0.050 & 0.093 & 0.482 & 0.029 & 0.091 & 583 \\
    DoS               & 0.608 & 0.532 & 0.117 & 0.149 & 0.196 & 0.233 & 4,089 \\
    Exploits          & 0.625 & 0.788 & 0.777 & 0.638 & 0.693 & 0.705 & 11,132 \\
    Fuzzers           & 0.293 & 0.275 & 0.592 & 0.627 & 0.392 & 0.382 & 6,062 \\
    Generic           & 0.999 & 0.999 & 0.971 & 0.969 & 0.984 & 0.984 & 18,871 \\
    Normal            & 0.967 & 0.978 & 0.755 & 0.690 & 0.848 & 0.809 & 37,000 \\
    Reconnaissance    & 0.931 & 0.887 & 0.803 & 0.821 & 0.863 & 0.853 & 3,496 \\
    Shellcode         & 0.342 & 0.187 & 0.701 & 0.878 & 0.460 & 0.308 & 378 \\
    Worms$^\dagger$   & 0.438 & 0.366 & 0.159 & 0.591 & 0.233 & 0.452 & 44 \\
    \midrule
    \textbf{Macro avg}    & 0.522 & 0.510 & 0.498 & 0.601 & 0.470 & 0.487 & 82,332 \\
    \textbf{Weighted avg} & 0.826 & 0.836 & 0.754 & 0.716 & 0.778 & 0.763 & 82,332 \\
    \bottomrule
    \multicolumn{8}{l}{\footnotesize B = Baseline; S = SMOTE; P = Precision;
      R = Recall. Values rounded to 3~d.p.}\\
    \multicolumn{8}{l}{\footnotesize $^\dagger$ $n=44$; metrics are indicative
      only at this sample size.}\\
  \end{tabular}
\end{table*}

% --- TABLE II ----------------------------------------------
\begin{table}[tb]
  \caption{Bootstrap 95\% CI for Aggregate Metrics
           ($n_{\text{resample}}=2{,}000$, Percentile)}
  \label{tab:bootstrap}
  \centering
  \begin{tabular}{llrrrr}
    \toprule
    \textbf{Model} & \textbf{Metric} &
      \textbf{Obs.} & \textbf{CI$_\text{lo}$} &
      \textbf{CI$_\text{hi}$} & \textbf{Width} \\
    \midrule
    Baseline & Accuracy    & 0.7543 & 0.7516 & 0.7574 & 0.0057 \\
    Baseline & Macro F1    & 0.4704 & 0.4556 & 0.4853 & 0.0298 \\
    Baseline & Weighted F1 & 0.7780 & 0.7754 & 0.7808 & 0.0054 \\
    SMOTE    & Accuracy    & 0.7161 & 0.7129 & 0.7191 & 0.0062 \\
    SMOTE    & Macro F1    & 0.4874 & 0.4752 & 0.4989 & 0.0237 \\
    SMOTE    & Weighted F1 & 0.7631 & 0.7605 & 0.7659 & 0.0054 \\
    \bottomrule
    \multicolumn{6}{l}{\footnotesize Non-overlapping accuracy CIs corroborate
      McNemar's test.}\\
  \end{tabular}
\end{table}

% ----------------------------------------------------------
\subsection{Explanation Quality (RQ2)}

\begin{figure}[t]
  \centering
  \includegraphics[width=0.49\columnwidth]{shap_bar_baseline.png}\hfill
  \includegraphics[width=0.49\columnwidth]{shap_bar_smote.png}
  \caption{\SHAP global feature importance (mean $|\text{SHAP}|$, 60 instances).
           (left)~Baseline, top: \texttt{sbytes} (0.0185).
           (right)~\SMOTE, top: \texttt{ct\_dst\_src\_ltm} (0.0204).
           Spearman $\rho=0.900$; top-5 overlap $=0.80$.}
  \label{fig:shap}
\end{figure}

\begin{figure}[t]
  \centering
  \includegraphics[width=0.49\columnwidth]{lime_importance_baseline.png}\hfill
  \includegraphics[width=0.49\columnwidth]{lime_importance_smote.png}
  \caption{\LIME global feature importance (mean $|\text{weight}|$, 60 instances).
           (left)~Baseline, top: \texttt{sttl} (0.042); $R^2=0.290$.
           (right)~\SMOTE, top: \texttt{service} (0.028); $R^2=0.252$.
           Spearman $\rho=0.565$; top-5 overlap $=0.60$.}
  \label{fig:lime}
\end{figure}

\SHAP global rankings showed moderate stability (Spearman $\rho=0.900$, top-5
overlap $=0.80$) (Fig.~\ref{fig:shap}). Wilcoxon test on 42 paired importances
confirmed a statistically significant shift ($W=283$, $p=0.0347$; Holm
$\alpha=0.050$) with a medium effect ($r=0.3259$) (Table~\ref{tab:hyp}).

\LIME rankings showed considerably lower stability ($\rho=0.565$, overlap $=0.60$)
(Fig.~\ref{fig:lime}). Wilcoxon on 36 paired \LIME importances confirmed a larger
shift ($W=115$, $p=3.59\times10^{-4}$; Holm $\alpha=0.025$) with a large effect
($r=0.5947$) (Table~\ref{tab:hyp}).

\LIME local surrogate fidelity, measured by the mean coefficient of determination
($R^2$) of each instance's local linear model, was 0.290 for the Baseline model
and 0.252 for the SMOTE model. These values indicate that the \LIME surrogate
captures approximately 25--29\% of local decision variance. The observed ranking
shift ($r=0.5947$) is statistically robust; however, \LIME attribution magnitudes
reflect the linear surrogate under moderate fidelity conditions and should be
understood as approximations of local model behaviour rather than exact
attributions. Inter-method agreement was low but improved slightly after
\SMOTE ($\rho=0.410\to0.476$).

\smallskip
\noindent\textbf{Summary (RQ2):} \SMOTE caused significant \XAI explanation shifts;
\LIME exhibited a large effect vs.\ medium for \SHAP. \LIME results interpreted
in light of moderate local surrogate fidelity ($R^2 \approx 0.25$--$0.29$).

% ----------------------------------------------------------
\subsection{Predictive vs.\ Explanatory Sensitivity (RQ3)}

\begin{figure*}[t]
  \centering
  \includegraphics[width=\textwidth]{explanation_ranking_comparison.png}
  \caption{Feature rank shift: Baseline vs.\ \SMOTE for \SHAP (top) and \LIME
           (bottom). Crossing lines indicate rank changes; \LIME shows
           substantially greater rank instability. Note: the two metrics operate
           on different scales and sample sizes; the figure illustrates qualitative
           asymmetry, not a precise proportional ratio.}
  \label{fig:ranking}
\end{figure*}

The central finding is a systematic asymmetry between predictive and explanatory
sensitivity (Fig.~\ref{fig:ranking}, Table~\ref{tab:effect}). All predictive
effect sizes are negligible (max $|h|=0.0868$), while explanation effect sizes
range from small (confidence scores, $r=0.1686$) to medium (\SHAP, $r=0.3259$)
to large (\LIME, $r=0.5947$). As these metrics operate on different scales and
reflect different sample sizes ($n=36$--42 feature pairs vs.\ $n=82{,}332$
paired predictions), the comparison is qualitative: explanation attribution
rankings are substantially more sensitive to class rebalancing than aggregate
classification outcomes.

\smallskip
\noindent\textbf{Summary (RQ3):} Explanatory sensitivity substantially exceeds
predictive sensitivity; \LIME attribution rankings exhibited the greatest shift
of all outputs examined.

% --- TABLE III ---------------------------------------------
\begin{table}[tb]
  \caption{Effect Sizes: Baseline vs.\ \SMOTE}
  \label{tab:effect}
  \centering
  \setlength{\tabcolsep}{4pt}
  \renewcommand{\arraystretch}{0.75}
  \begin{tabular}{llrl}
    \toprule
    \textbf{Comparison} & \textbf{Metric} & \textbf{Value} & \textbf{Magnitude} \\
    \midrule
    Accuracy         & $h$ & $-0.0868$ & Negligible \\
    Macro F1         & $h$ & $+0.0340$ & Negligible \\
    Weighted F1      & $h$ & $-0.0354$ & Negligible \\
    Confidence shift & $r$ & $0.1686$  & Small \\
    \SHAP shift      & $r$ & $0.3259$  & Medium \\
    \LIME shift      & $r$ & $0.5947$  & Large \\
    \bottomrule
    \multicolumn{4}{p{\linewidth}}{\footnotesize $h$: Eq.~\eqref{eq:cohenh} ($n=82{,}332$);
      $r$: Eq.~\eqref{eq:rbs} ($n=42$ \SHAP, $n=36$ \LIME).
      Thresholds (Cohen~1988): $|h|{<}0.20$ neg.; $r{\geq}0.10$ sm.,
      ${\geq}0.30$ med., ${\geq}0.50$ lg. Cross-metric: qualitative only.}\\
  \end{tabular}
\end{table}

% ----------------------------------------------------------
\subsection{Hypothesis Test Summary (RQ4)}

All four tests rejected $H_0$ after Holm--Bonferroni correction (Table~\ref{tab:hyp}).
The asymmetry between negligible predictive effect sizes and medium-to-large
explanation effect sizes supports the research hypothesis. \LIME's greater
sensitivity than \SHAP ($\Delta r=0.269$) is consistent with the expectation that
perturbation-based local methods are more responsive to training distribution
changes than exact tree-attribution methods.

% --- TABLE IV ----------------------------------------------
\begin{table}[tb]
  \caption{Hypothesis Test Summary (Holm--Bonferroni, $\alpha=0.05$)}
  \label{tab:hyp}
  \centering
  \footnotesize
  \setlength{\tabcolsep}{3.5pt}
  \resizebox{\columnwidth}{!}{%
  \begin{tabular}{lllrl}
    \toprule
    \textbf{Test} & \textbf{Statistic} &
      $\bm{p}$\textbf{-value} & \textbf{Holm~}$\bm{\alpha}$ & \textbf{Reject} \\
    \midrule
    McNemar (Yates)  & $\chi^2=1547.51$    & $<10^{-300}$        & 0.0125 & Yes \\
    Wilcoxon (conf.) & $W=455{,}648{,}244$ & $<10^{-300}$        & 0.0167 & Yes \\
    Wilcoxon (\LIME) & $W=115.0$           & $3.59\times10^{-4}$ & 0.0250 & Yes \\
    Wilcoxon (\SHAP) & $W=283.0$           & $3.47\times10^{-2}$ & 0.0500 & Yes \\
    \bottomrule
    \multicolumn{5}{p{\linewidth}}{\footnotesize Ordered by Holm rank. $b=1{,}632$,
      $c=4{,}784$ (McNemar). \SHAP power $\approx0.55$ ($n=42$);
      \LIME power $\approx0.90$ ($n=36$).}\\
  \end{tabular}}
\end{table}

% ==========================================================
\section{Discussion}
\label{sec:discussion}
% ==========================================================

\subsection{Interpretation of Findings}

The central result --- that \SMOTE changes explanation structure more than
prediction structure --- follows from a mechanistic argument. Class rebalancing
redistributes the model's learned decision boundary: minority classes receive
proportionally more training signal, shifting internal feature weighting even when
aggregate accuracy is little changed. Post-hoc explanation methods detect this
shift directly because they query the local decision surface rather than summarising
it with a single metric. \LIME is more sensitive than \SHAP because it evaluates
the model on perturbed local neighbourhoods directly affected by the shifted
boundary, whereas \SHAP's global Shapley-value attributions partially average out
local perturbations.

\subsection{Practical Implications}

\noindent\textbf{(1) Retraining validation:} Accuracy alone is insufficient.
\XAI outputs may change substantially even when accuracy is stable; explanation
outputs must be re-validated after any retraining event. \textbf{(2) Deployment
trade-off:} Of 82,332 test instances, \SMOTE gained 1,632 predictions (1.98\%)
but lost 4,784 (5.81\%) vs.\ Baseline; for rare-attack-priority deployments this
is acceptable, but for false-positive-sensitive environments it warrants
consideration. \textbf{(3) Method
selection:} A distribution change can produce a large shift ($r=0.595$) in \LIME
rankings without a corresponding accuracy change; \LIME-based pipelines require
independent validation after retraining. \textbf{(4) Operational monitoring:}
Tracking Spearman~$\rho$ of feature rankings before and after retraining detects
explanation drift that accuracy-only monitoring misses.

\vspace{-0.5em}
\subsection{Relationship to Trustworthy AI}

The predictive-explanatory asymmetry has direct implications for trustworthy-AI
frameworks in cybersecurity \cite{charmet2022,rjoub2023}. Regulatory guidance
increasingly requires stable explanations for consequential decisions
\cite{euaiact2024}; explanations that are unstable under routine model updates
become unreliable even when performance metrics remain acceptable. To the best of
the author's knowledge, this study provides the first empirical evidence of this
instability in a standardised \NIDS benchmark setting.

% ==========================================================
\section{Limitations}
\label{sec:limitations}
% ==========================================================

This study has the following limitations.

\textit{Internal validity}: Only \SMOTE was evaluated; ADASYN and cost-sensitive
methods may produce different explanatory outcomes. Results are specific to Random
Forest; \SHAP TreeExplainer is tree-specific and \LIME sensitivity may differ for
neural networks. Explanation subset size and statistical power: At $n=42$, \SHAP
Wilcoxon power $\approx0.55$ (below the conventional 0.80 threshold); the \SHAP
result ($p=0.0347$) should be interpreted with caution. At $n=36$, \LIME power
$\approx0.90$. The non-significant SHAP result would have been inconclusive at
conventional power thresholds.

\textit{Construct validity}: Explanation quality is operationalised as rank
stability and magnitude shift; human interpretability and ground-truth fidelity
are not assessed. LIME surrogate fidelity: The \LIME surrogate fidelity
($R^2 = 0.290$ Baseline, $R^2 = 0.252$ SMOTE) is moderate, meaning the linear
surrogate captures approximately 25--29\% of local decision variance. The ranking
shift is statistically robust, but absolute \LIME magnitudes are approximations.
SHAP vs.\ LIME comparability: The absolute magnitudes of their rank-biserial $r$
values are not directly commensurate. \SHAP operates on 42 paired feature rankings
($n=42$); \LIME operates on 36 paired feature rankings ($n=36$). Cross-method
comparison indicates relative sensitivity but cannot establish proportional
relationships. Worms class sample size: The Worms test class contains only 44
instances. F1 comparisons at this scale are sensitive to individual prediction
outcomes and are reported as indicative results rather than stable population
estimates.

\textit{External validity}: UNSW-NB15 is a point-in-time capture from 2015.
Relevant attack techniques have evolved significantly: TLS~1.3 encryption,
microservice traffic patterns, and cloud APIs create operational conditions
substantially different from 2015 laboratory traffic. Transferability to modern
encrypted-traffic environments warrants validation on CICIDS-2017 or
CSE-CIC-IDS-2018. Concept drift in live network traffic may invalidate both
predictive performance and explanation patterns observed in this study.

% ==========================================================
\section{Future Work}
\label{sec:future}
% ==========================================================

Future work should address: (i)~alternative rebalancing techniques (ADASYN,
class-weighted loss); (ii)~gradient-based explanations for neural network \NIDS;
(iii)~multi-dataset validation (CICIDS-2017, NSL-KDD); (iv)~human-centred
evaluation of explanation rank shifts on analyst decision-making; and
(v)~formal explanation-consistency metrics for model certification.

% ==========================================================
\IEEEtriggerfinalcols
\section{Conclusion}
\label{sec:conclusion}
% ==========================================================

A controlled empirical study on UNSW-NB15 shows that \SMOTE-based class
rebalancing produces statistically significant but practically negligible changes
to aggregate \NIDS predictive metrics (max $|h|=0.087$), while producing
medium-to-large changes to \SHAP and \LIME feature attribution rankings
(rank-biserial $r$ up to 0.595 for \LIME). \LIME's explanatory effect size
substantially and disproportionately exceeds predictive sensitivity across
complementary effect-size measures. The disproportionate sensitivity of
explanation outputs relative to classification outcomes demonstrates that
accuracy-only validation is insufficient for \XAI-augmented \NIDS pipelines.

The secondary finding --- that \LIME is substantially more sensitive than \SHAP
to training distribution changes ($\Delta r=0.269$) --- has direct implications
for \XAI tool selection: perturbation-based local methods must be independently
validated when training data composition changes, and surrogate fidelity should
be reported alongside attribution rankings. This practice enables practitioners to
distinguish ranking shifts caused by model changes from shifts caused by surrogate
approximation variance. This finding holds even with the caveat that \LIME local
surrogate fidelity ($R^2 \approx 0.25$--$0.29$) is moderate: the statistical
evidence for a large explanatory shift is robust regardless of absolute attribution
magnitudes.

These results support explanation-consistency monitoring as an important component
of responsible model assessment in any \NIDS pipeline that incorporates post-hoc
\XAI methods, particularly in regulatory contexts that require stable and reliable
explanations \cite{euaiact2024}.

% ==========================================================
%  REFERENCES  (IEEE appearance order, 23 entries)
% ==========================================================
\begin{thebibliography}{23}

\bibitem{alshamy2021}
R.~Alshamy, M.~Ghurab, S.~Othman, and F.~Alshami,
``Intrusion detection model for imbalanced dataset using SMOTE and random
forest algorithm,''
in \emph{Advances in Cyber Security (ACeS~2021)},
ser. Commun. Comput. Inf. Sci. (CCIS), vol.~1487.
Singapore: Springer, 2021, pp.~311--323.
doi:~10.1007/978-981-16-8059-5\_22.

\bibitem{wu2022}
T.~Wu, H.~Fan, H.~Zhu, C.~You, H.~Zhou, and X.~Huang,
``Intrusion detection system combined enhanced random forest with SMOTE
algorithm,''
\emph{EURASIP J. Adv. Signal Process.}, vol.~2022, Art. no.~39, 2022.
doi:~10.1186/s13634-022-00871-6.

\bibitem{more2024}
S.~More, M.~Idrissi, H.~Mahmoud, and A.~T.~Asyhari,
``Enhanced intrusion detection systems performance with UNSW-NB15 data
analysis,''
\emph{Algorithms}, vol.~17, no.~2, Art. no.~64, Feb. 2024.
doi:~10.3390/a17020064.

\bibitem{moustafa2015}
N.~Moustafa and J.~Slay,
``UNSW-NB15: A comprehensive data set for network intrusion detection
systems,''
in \emph{Proc. IEEE Military Commun. Inf. Syst. Conf. (MilCIS)},
Canberra, Australia, Nov. 2015, pp.~1--6.
doi:~10.1109/MilCIS.2015.7348942.

\bibitem{shanmugam2025}
V.~Shanmugam, R.~Razavi-Far, and E.~Hallaji,
``Addressing class imbalance in intrusion detection: A comprehensive
evaluation of machine learning approaches,''
\emph{Electronics}, vol.~14, no.~1, Art. no.~69, Jan. 2025.
doi:~10.3390/electronics14010069.

\bibitem{chawla2002}
N.~V.~Chawla, K.~W.~Bowyer, L.~O.~Hall, and W.~P.~Kegelmeyer,
``SMOTE: Synthetic minority over-sampling technique,''
\emph{J. Artif. Intell. Res.}, vol.~16, pp.~321--357, Jun. 2002.
doi:~10.1613/jair.953.

\bibitem{sayegh2024}
H.~R.~Sayegh, W.~Dong, and A.~M.~Al-madani,
``Enhanced intrusion detection with LSTM-based model, feature selection,
and SMOTE for imbalanced data,''
\emph{Appl. Sci.}, vol.~14, no.~2, Art. no.~479, Jan. 2024.
doi:~10.3390/app14020479.

\bibitem{charmet2022}
F.~Charmet, H.~C.~Tanuwidjaja, S.~Ayoubi, P.~F.~Gimenez, Y.~Han,
H.~Jmila \emph{et~al.},
``Explainable artificial intelligence for cybersecurity: A literature
survey,''
\emph{Ann. Telecommun.}, vol.~77, no.~11--12, pp.~789--812, 2022.
doi:~10.1007/s12243-022-00926-7.

\bibitem{rjoub2023}
G.~Rjoub, J.~Bentahar, O.~Abdel~Wahab, R.~Mizouni, A.~Song,
R.~Cohen \emph{et~al.},
``A survey on explainable artificial intelligence for cybersecurity,''
\emph{IEEE Trans. Netw. Serv. Manag.}, vol.~20, no.~4,
pp.~5115--5140, Dec. 2023.
doi:~10.1109/TNSM.2023.3282740.

\bibitem{lundberg2017}
S.~M.~Lundberg and S.-I.~Lee,
``A unified approach to interpreting model predictions,''
in \emph{Adv. Neural Inf. Process. Syst. (NeurIPS)}, vol.~30,
Long Beach, CA, USA, Dec. 2017, pp.~4766--4777. [Online]. Available:
\url{https://proceedings.neurips.cc/paper_files/paper/2017/hash/%
8a20a8621978632d76c43dfd28b67767-Abstract.html}

\bibitem{lundberg2020}
S.~M.~Lundberg, G.~Erion, H.~Chen, A.~DeGrave, J.~M.~Prutkin,
B.~Nair \emph{et~al.},
``From local explanations to global understanding with explainable AI
for trees,''
\emph{Nat. Mach. Intell.}, vol.~2, no.~1, pp.~56--67, Jan. 2020.
doi:~10.1038/s42256-019-0138-9.

\bibitem{ribeiro2016}
M.~T.~Ribeiro, S.~Singh, and C.~Guestrin,
```Why should I trust you?': Explaining the predictions of any
classifier,''
in \emph{Proc. 22nd ACM SIGKDD Int. Conf. Knowl. Discovery Data Mining
(KDD)}, San Francisco, CA, USA, Aug. 2016, pp.~1135--1144.
doi:~10.1145/2939672.2939778.

\bibitem{gaspar2024}
D.~Gaspar, P.~Silva, and C.~Silva,
``Explainable AI for intrusion detection systems: LIME and SHAP
applicability on multi-layer perceptron,''
\emph{IEEE Access}, vol.~12, pp.~30164--30175, 2024.
doi:~10.1109/ACCESS.2024.3368377.

\bibitem{hermosilla2025a}
P.~Hermosilla, S.~Berr\'ios, and H.~Allende-Cid,
``Explainable AI for forensic analysis: A comparative study of SHAP and
LIME in intrusion detection models,''
\emph{Appl. Sci.}, vol.~15, no.~13, Art. no.~7329, Jul. 2025.
doi:~10.3390/app15137329.

\bibitem{hermosilla2025b}
P.~Hermosilla, M.~D\'iaz, S.~Berr\'ios, and H.~Allende-Cid,
``Use of explainable artificial intelligence for analyzing and explaining
intrusion detection systems,''
\emph{Computers}, vol.~14, no.~5, Art. no.~160, May 2025.
doi:~10.3390/computers14050160.

\bibitem{visani2022}
G.~Visani, E.~Bagli, F.~Chesani, A.~Poluzzi, and D.~Capuzzo,
``Statistical stability indices for LIME: Obtaining reliable explanations
for machine learning models,''
\emph{J. Oper. Res. Soc.}, vol.~73, no.~1, pp.~91--101, 2022.
doi:~10.1080/01605682.2020.1865846.

\bibitem{patil2020}
A.~Patil, A.~Framewala, and F.~Kazi,
``Explainability of SMOTE based oversampling for imbalanced dataset
problems,''
in \emph{Proc. IEEE Int. Conf. Inf. Commun. Technol. (ICICT)},
Feb. 2020, pp.~41--45.
doi:~10.1109/ICICT50521.2020.9092325.

\bibitem{alarab2022}
I.~Alarab and S.~Prakoonwit,
``Effect of data resampling on feature importance in imbalanced blockchain
data: Comparison studies of resampling techniques,''
\emph{Data Sci. Manag.}, vol.~5, no.~2, pp.~66--76, Jun. 2022.
doi:~10.1016/j.dsm.2022.04.003.

\bibitem{breiman2001}
L.~Breiman,
``Random forests,''
\emph{Mach. Learn.}, vol.~45, no.~1, pp.~5--32, Oct. 2001.
doi:~10.1023/A:1010933404324.

\bibitem{virtanen2020}
P.~Virtanen, R.~Gommers, T.~E.~Oliphant, M.~Haberland, T.~Reddy,
D.~Cournapeau \emph{et~al.},
``SciPy~1.0: Fundamental algorithms for scientific computing in Python,''
\emph{Nat. Methods}, vol.~17, pp.~261--272, Feb. 2020.
doi:~10.1038/s41592-020-0772-5.

\bibitem{pedregosa2011}
F.~Pedregosa, G.~Varoquaux, A.~Gramfort, V.~Michel, B.~Thirion,
O.~Grisel \emph{et~al.},
``Scikit-learn: Machine learning in Python,''
\emph{J. Mach. Learn. Res.}, vol.~12, pp.~2825--2830, Nov. 2011.
[Online]. Available: \url{http://jmlr.org/papers/v12/pedregosa11a.html}

\bibitem{lemaitre2017}
G.~Lema\^{i}tre, F.~Nogueira, and C.~K.~Aridas,
``Imbalanced-learn: A Python toolbox to tackle the curse of imbalanced
datasets in machine learning,''
\emph{J. Mach. Learn. Res.}, vol.~18, no.~17, pp.~1--5, 2017.
[Online]. Available: \url{http://jmlr.org/papers/v18/16-365.html}

\bibitem{euaiact2024}
European Parliament and Council of the European Union,
``Regulation (EU) 2024/1689 laying down harmonised rules on artificial
intelligence (Artificial Intelligence Act),''
\emph{Off. J. Eur. Union}, Jul. 2024. [Online]. Available:
\url{https://eur-lex.europa.eu/legal-content/EN/TXT/?uri=CELEX:32024R1689}

\end{thebibliography}

% ==========================================================
\end{document}
% ==========================================================
